\section{Structural Negatives}\label{sec:structural-negatives}

We collect here a second family of rigorous no-go results.  Unlike the proof
class barriers of \Cref{sec:proof-class-barriers}, these statements are best
read as structural diagnoses of particular candidate mechanisms.  Each one says
that a natural route fails for a concrete mathematical reason.

\subsection{Universal Block-Product Carrier-Mass Counter}

% source: phase4/block_product_carrier_mass_note.md
\begin{proposition}\label{prop:block-product-counter}
For every Shortener strategy there exists a Prolonger response $P$ such that
\[
\sum_{p\in B(P)}\frac1p \ge \frac12\log\log n + O(1),
\]
where $B(P)$ is the set of prime carriers used by Prolonger's moves.
Consequently no proof of $\L(n)=O(n/\log n)$ can proceed by first showing that a
Shortener strategy forces uniformly bounded reciprocal carrier mass.
\end{proposition}

\begin{proof}[Proof sketch]
Partition the primes $p\le n^{1/3}$ into disjoint blocks whose products lie in
$(\sqrt n,n^{5/6}]$.  These block products are pairwise coprime, so a single
Shortener move can kill at most one of them.  Greedy descending-by-mass play
then lets Prolonger secure at least half of the total reciprocal mass of the
small primes, namely $(1/2)\log\log n+O(1)$.  The standard smooth-number sieve
would convert a uniform $O(1)$ carrier-mass bound into $O(n/\log n)$, so the
carrier-mass route is dead in that form.
\end{proof}

\subsection{Separate-Rank Fan No-Go}

% source: researcher-17-codex-negative-on-separate-rank-fan-route.md
\begin{proposition}\label{prop:separate-rank-fan}
Let $H(n)=\sum_{p\le n^\delta}1/p\sim\log\log n$.  In the fixed-rank upper-half
fan architecture,
\[
W_h(n)\asymp \frac{n}{\log n}\frac{H(n)^h}{h!},
\]
so
\[
\max_h W_h(n)\asymp \frac{n}{\sqrt{\log\log n}}=o(n).
\]
With the proved divisor-shadow retention factor
$2^{-2(2^h-1)}$, the total separate-rank contribution satisfies
\[
\sum_{h\ge 1} F_h(n)
=
\frac{n}{(\log n)^{1-o(1)}}=o(n).
\]
Hence neither a single rank nor a rank-by-rank summation can yield a linear
lower bound.
\end{proposition}

\begin{proof}[Proof sketch]
The single-rank ceiling follows by optimizing $H^h/h!$ at $h\asymp H$ and
applying Stirling.  For the full route, one maximizes
\[
F_h(n)=\frac{n}{\log n}\frac{H^h}{h!}\,2^{-2(2^h-1)}
\]
over $h$; the optimum occurs at $h\asymp \log_2\log\log\log n$, and even there
the gain is only $(\log n)^{o(1)}$.  Summing over all ranks changes only the
$o(1)$ term in the exponent.
\end{proof}

\subsection{Multi-Rank Same-$b$ Coupling Obstruction}

% source: current_state.md (Round 18 closeout)
\begin{proposition}\label{prop:same-b-coupling}
The obvious same-$b$ multi-rank coupling mechanism for the fan hierarchy is
blocked twice:
\begin{enumerate}
  \item same-$b$ upper-half cores cannot be nested (Sperner obstruction);
  \item each earlier compatible move shields at most one top-lateral divisor of
  a given upper-half target (LCM obstruction).
\end{enumerate}
Consequently the hoped-for uniform multi-rank shielding theorem cannot be
derived from same-$b$ coupling.
\end{proposition}

\begin{proof}[Proof sketch]
If $S\subsetneq T$ and both $A_S b$ and $A_T b$ lie in $(n/2,n]$, then
$A_Tb=A_Sb\cdot A_{T\setminus S}\ge 2A_Sb>n$, impossible.  So the same-$b$
cores form a Sperner family rather than a nested chain.  For the LCM
obstruction, if one earlier move $z\le n$ shielded two top-lateral divisors of
$y=A_Tb$, then the least common multiple of those divisors would equal $y$,
forcing $y\mid z$ and hence $z=y$, contradicting the requirement that $z$ be an
earlier distinct move.
\end{proof}

\subsection{Residual-Width Refutation}

% source: researcher-20-codex-residual-width-framework-refuted.md
\begin{proposition}\label{prop:residual-width}
The domination-only residual-width lemma is false.  More precisely, for every
certificate family $C\subseteq \{2,\dots,n\}$ there exists an antichain
$A\subseteq \U$ with $|A|\le |C|$ such that $C\subseteq \operatorname{Comp}(A)$ and
\[
w(R(A))\ge \frac n2-|C|-O(1).
\]
Thus residual width cannot be controlled from $|C|$ and domination data alone.
\end{proposition}

\begin{proof}[Proof sketch]
For $c\in C$, let $\lambda_n(c)$ be the least multiple of $c$ in $(n/2,n]$.
Then $A=\{\lambda_n(c):c\in C\}$ is an antichain in $\U$, each $c\in C$ is
comparable with some element of $A$, and every untouched element of $\U$ is
incomparable with all of $A$.  Hence $\U\setminus A\subseteq R(A)$, which
already gives the lower bound on the residual width.
\end{proof}

\subsection{Ford-Route Correction}

% source: researcher-23-codex-ford-route-correction.md
\begin{proposition}\label{prop:ford-correction}
The old Ford-band reduction based only on a factorization
$u=am$ with $a\in (Y,2Y]$ and $P^-(m)>Y^\delta$ is invalid.  Roughness of the
cofactor is not enough: Ford's divisor-interval theorem applies to integers
that are themselves rough, and the naive bandwise conclusion can fail by a full
factor of $\log n$.
\end{proposition}

\begin{proof}[Proof sketch]
The counterexample takes a dyadic band of smooth squarefree skeletons
$a=30c\in (Y,2Y]$ with constant reciprocal mass and multiplies them by large
primes $\ell\in (n/(2a),n/a]$.  Every resulting upper-half integer has a
factorization $u=a\ell$ with rough cofactor $\ell$, but there are
$\gg n/\log n$ such numbers in one band, not $O(n/\log^2 n)$.  So the reduction
was importing the wrong theorem; any valid Ford-based route must control the
roughness of the whole survivor or add separate bandwise reciprocal-mass
control.
\end{proof}

\subsection{Band-Local Closure Explosion}

% source: researcher-24-codex-band-local-closure-explosion.md
\begin{proposition}\label{prop:band-local-explosion}
Fix $\lambda>1$ and $0<\gamma<1$.  Let
\[
\mathcal P_{\mathrm{hi}}=\{p\text{ prime}:X\le p\le X^\lambda\},
\qquad
\mathcal C_X(\gamma)=\{c\le X^\gamma:\mu^2(c)=1\},
\]
and let $\mathcal B_X(\lambda)$ be the set of triple products $pqr$ with
$p<q<r$ in $\mathcal P_{\mathrm{hi}}$.  Then some dyadic band $(Y,2Y]$ contains
a set
\[
\mathcal A_Y\subseteq \{bc:b\in\mathcal B_X(\lambda),\ c\in\mathcal C_X(\gamma)\}
\]
with
\[
\sum_{a\in\mathcal A_Y}\frac1a\gg_{\lambda,\gamma} 1.
\]
If the proper divisors of the elements of $\mathcal A_Y$ are already
unavailable and $Y\le n^{1-\eta}$, this single band supports
$\gg_{\lambda,\gamma,\eta} n/\log n$ legal upper-half moves.
\end{proposition}

\begin{proof}[Proof sketch]
The reciprocal mass of the triple products in the high band is constant:
\[
\sum_{p<q<r\in \mathcal P_{\mathrm{hi}}}\frac1{pqr}\gg_\lambda 1.
\]
The reciprocal mass of the low squarefree pool is logarithmic:
\[
\sum_{c\in\mathcal C_X(\gamma)}\frac1c\asymp_\gamma \log X.
\]
Their product therefore has total reciprocal mass $\gg \log X$, and pigeonholing
over only $O(\log X)$ dyadic bands yields one band with constant mass.  The
upper-half lift then follows from the prime number theorem in the intervals
$(n/(2a),n/a]$.
\end{proof}

\subsection{Directed Rank-3 Cleanup Is Not Budget-Limited}

% source: researcher-25-codex-directed-rank3-budget.md
\begin{proposition}\label{prop:rank3-budget}
For a fixed-power band $I=(X,X^\lambda]\subseteq [2,n^\alpha]$ with
$1/3<\alpha<1/2$, let
\[
T_I(n)=\#\{p<q<r:\ p,q,r\in I,\ pqr\le n/2\}.
\]
Then
\[
T_I(n)\ll_{\alpha,\lambda}\frac{n}{\log n},
\qquad
\sum_I T_I(n)\ll_{\alpha,\lambda}\frac{n\log\log n}{\log n}.
\]
Thus the natural same-band rank-$3$ cleanup is not blocked by raw move count;
the real obstruction is online fiber persistence after Prolonger preempts one
triple.
\end{proposition}

\begin{proof}[Proof sketch]
For fixed $p,q\in I$, the admissible primes $r$ satisfy
$r\le n/(2pq)$, so the prime number theorem bounds their number by
$\ll n/(pq\log n)$.  Summing over $p<q$ in one fixed-power band and using
Mertens on that band gives $T_I(n)\ll n/\log n$.  Since there are only
$O(\log\log n)$ such bands up to $n^\alpha$, the total budget is
$O(n\log\log n/\log n)$, well below the conditional T2 scale.
\end{proof}

\subsection{Dyadic-Fiber Positive-Density Collapse}

% source: researcher-22-pro-dyadic-fiber-collapse.md
\begin{proposition}\label{prop:dyadic-fiber-collapse}
A positive-density theorem for one dyadic fiber is equivalent in strength to a
linear lower bound for the original game.  More precisely, if Prolonger could
force $\eta K$ moves inside one fiber
\[
\mathcal T_b=\{bc:K/2<c\le K\},
\qquad
b>\sqrt n,
\]
then $L(K)\gg_\eta K$ for the original divisibility game at scale $K$.
\end{proposition}

\begin{proof}[Proof sketch]
Inside one fixed $b$-fiber, divisibility is exactly divisibility in the scaled
game on $[2,K]$, because
\[
bd_1\mid bd_2 \Longleftrightarrow d_1\mid d_2.
\]
Moreover, different large-prime fibers are laterally incomparable when
$b>\sqrt n$.  So a Shortener strategy can play independently inside each fiber,
and a positive-density lower bound in one fiber would already imply a positive
density lower bound for the original upper-half divisibility game at scale $K$.
In other words, the would-be auxiliary theorem is just the main linear
conjecture in disguise.
\end{proof}
